\documentclass[aspectratio=169]{beamer}
\usepackage{etex} % fixes new-dimension error
\usepackage{lmodern}
\usepackage[T1]{fontenc}

\usepackage{graphicx,amsmath}
\usepackage{stmaryrd} % cf. interleave
\usepackage{booktabs}
\usepackage{amscd}
\usepackage{multicol}
\usepackage[absolute,overlay]{textpos}
\usepackage{alltt}
\usepackage{proof}
\usepackage{subcaption}
\usepackage{tikz}
\usepackage{tikz-cd}
\usepackage[new]{old-arrows}
\usepackage[all]{xy}
\usepackage{pgfplots}
\usepackage{textcomp}

\usepackage{transparent}
\usepackage{xspace}
\usepackage{listings}
\usepackage{pdfpages}
\usepackage{relsize}



%%%%%%%%%%%%% Macros
% \newcommand{\Ban}{\catfont{Ban}}
% \newcommand{\Cats}{\catfont{Cat}}

%%%% Misc
%% Operations
\newcommand{\sem}[1]{\llbracket #1 \rrbracket}
\DeclareMathOperator{\img}{\mathrm{im}}
\DeclareMathOperator{\dom}{\mathrm{dom}}
\DeclareMathOperator{\codom}{\mathrm{codom}}
%% Sets of numbers
\newcommand{\N}{\mathbb{N}}
\newcommand{\Z}{\mathbb{Z}}
\newcommand{\Nats}{\mathbb{N}}
\newcommand{\Reals}{\mathbb{R}}
\newcommand{\Rz}{\Reals_{\geq 0}}
\newcommand{\Complex}{\mathbb{C}}
%% Writing
\newcommand{\cf}{\emph{cf.}}
\newcommand{\ie}{\emph{i.e.}}
\newcommand{\eg}{\emph{e.g.}}
\newcommand{\df}[1]{\emph{\textbf{#1}}}
%%%%%%%%%%%%%%%% End of Misc

%%%% Programming Stuff
%% Types
\newcommand{\typefont}[1]{\mathbb{#1}}
\newcommand{\typeOne}{1}
\newcommand{\typeTwo}{2}
\newcommand{\typeA}{\typefont{A}}
\newcommand{\typeX}{\typefont{X}}
\newcommand{\typeB}{\typefont{B}}
\newcommand{\typeC}{\typefont{C}}
\newcommand{\typeV}{\typefont{V}}
\newcommand{\typeD}{\typefont{D}}
\newcommand{\typeI}{\typefont{I}}
%% RuleName
\newcommand{\rulename}[1]{(\mathrm{#1})}
%% Sequents
\newcommand{\jud}{\vdash}
\newcommand{\vljud}{\rhd}
\newcommand{\cojud}{\vdash_{\co}}
\newcommand{\vl}{\mathtt{v}}
\newcommand{\co}{\mathtt{c}}
% Program font
\newcommand{\prog}[1]{\mathtt{#1}}
\newcommand{\pseq}[3]{#1 \leftarrow #2; #3}
\newcommand{\ppm}[4]{(#1,#2) \leftarrow #3; #4}
\newcommand{\pinl}[1]{\prog{inl}(#1)}
\newcommand{\pinr}[1]{\prog{inr}(#1)}
\newcommand{\pcase}[4]{\prog{ case } #1 \prog{ of } \pinl{#2} \Rightarrow #3 ; \pinr{#2} \Rightarrow #4}
%% Sets of terms
\newcommand{\ValuesBP}[2]{\mathsf{Values}(#1, #2)}
\newcommand{\TermsBP}[2]{\mathsf{Terms}(#1, #2)}
\newcommand{\closValP}[1]{\ValuesBP{\emptyset}{#1}}
\newcommand{\closTermP}[1]{\TermsBP{\emptyset}{#1}}
\newcommand{\closVal}{\closValP{\typeA}}
\newcommand{\closTerm}{\closTermP{\typeA}}
%% Contextual equivalence
\newcommand{\ctxeq}{\equiv_{\prog{ctx}}}
%%%% End of Programming Stuff


% Misc by José
\newcommand{\wrap}[2][]{\begin{tabular}[#1]{@{}c@{}}#2\end{tabular}}
\newcommand{\mwrap}[1]{\ensuremath{\begin{array}{@{}c@{}}#1\end{array}}}
\def\trans#1{\xrightarrow{#1}}  % - a - > 
\def\Trans#1{\stackrel{#1}{\Longrightarrow}} % =a=> 
\newcommand{\transp}[2][35]{\color{fg!#1}#2}
\newcommand{\transpt}[2][.35]{\tikz{\node[inner sep=1pt,fill opacity=0.5]{#2}}}
\newcommand{\faded}[2][0.4]{{\transparent{#1}#2}} % alternative to "transp" using transparent package
\newcommand{\set}[1]{\left\{ #1 \right\}} % {a,b,...z}
\newcommand{\mi}[1]{\ensuremath{\mathit{#1}}\xspace}
\newcommand{\mf}[1]{\ensuremath{\mathsf{#1}}\xspace}
% \newcommand{\gold}[1]{\textcolor{darkgoldenrod}{#1}\xspace}


%------ using color ---------------------------------------------------------
\definecolor{goldenrod}{rgb}{.80392 .60784 .11373}
\definecolor{darkgoldenrod}{rgb}{.5451 .39608 .03137}
\definecolor{brown}{rgb}{.15 .15 .15}
\definecolor{darkolivegreen}{rgb}{.33333 .41961 .18431}
\definecolor{myGray}{gray}{0.85}
%
%
\newcommand{\red}[1]{\textcolor{red!80!black}{#1}\xspace}
\newcommand{\blue}[1]{\textcolor{blue}{#1}\xspace}
\newcommand{\gold}[1]{\textcolor{darkgoldenrod}{#1}\xspace}
\newcommand{\gray}[1]{\textcolor{myGray}{#1}\xspace}
% \def\alert#1{{\darkgoldenrod #1}}
% \def\alert#1{{\alert{#1}}}
%\def\brw#1{{\brown #1}}
% \def\structure#1{{\blue #1}}
% \def\tstructure#1{\textbf{\darkblue #1}}
%%\def\gre#1{{\green #1}}
\def\gre#1{{\darkolivegreen #1}}
\def\gry#1{{\textcolor{gray}{#1}}}
\def\rdb#1{{\red #1}}
\def\st{\mathbf{.}\,}
\def\laplace#1#2{*\txt{\mbox{ \fcolorbox{black}{myGray}{$\begin{array}{c}\mbox{#1}\\\\#2\\\\\end{array}$} }}}
%\newcommand{\galois}[2]{#1\; \dashv\; #2}



\def\ainv#1{\overline{#1}}
\def\aconv#1{#1^{\circ}} 
\def\rtran#1{\stackrel{#1}{\longrightarrow}}
\def\pair#1{\langle #1 \rangle}
\def\setdef#1#2{\mathopen{\{} #1 \asor #2 \mathclose{\}}}
\def\imp{\mathbin{\Rightarrow}}
\def\dimp{\mathbin{\Leftrightarrow}}
\def\rimp{\mathbin{\Leftarrow}}
\def\rra{\longrightarrow}
% \def\rcb#1#2#3#4{\def\nothing{}\def\range{#3}\mathopen{\langle}#1 \ #2 \ \ifx\range\nothing::\else: \ #3 :\fi \ #4\mathclose{\rangle}}
\def\abv{\stackrel{\rm abv}{=}}



% Spliting frames in 2 columns
\newcommand{\splittwo}[4]{ 
  \begin{columns}[T]% align columns
  \begin{column}{#1\textwidth} #3 \end{column} ~~~
  \begin{column}{#2\textwidth} #4 \end{column} \end{columns}
}\newcommand{\frsplit}[3][.48]{
  \begin{columns}%[T] % align columns
  \begin{column}{#1\textwidth} #2 \end{column} ~~~
  \begin{column}{#1\textwidth} #3 \end{column} \end{columns}
}
\newcommand{\frsplitdiff}[5][]{
  \begin{columns}[#1]%[T] % align columns
  \begin{column}{#2\textwidth} #4 \end{column} ~~~
  \begin{column}{#3\textwidth} #5 \end{column} \end{columns}
}
\newcommand{\frsplitt}[3][.48]{
  \begin{columns}[T] % align columns
  \begin{column}{#1\textwidth} #2 \end{column} ~~~
  \begin{column}{#1\textwidth} #3 \end{column} \end{columns}
}
\newcommand{\col}[2][.48]{\begin{column}{#1\textwidth} #2 \end{column}}
\newcommand{\colb}[3][.48]{\begin{column}{#1\textwidth} \begin{block}{#2} #3 \end{block} \end{column}}

% Spliting frames in 3 columns
\newcommand{\splitthree}[6]{
  \begin{columns}[T] % align columns
  \begin{column}{#1\textwidth} #4 \end{column} ~~~
  \begin{column}{#2\textwidth} #5 \end{column} ~~~
  \begin{column}{#3\textwidth} #6 \end{column} \end{columns}
}
\newcommand{\frsplitthree}[4][.31]{
  \begin{columns}%[T] % align columns
  \begin{column}{#1\textwidth} #2 \end{column} ~~~
  \begin{column}{#1\textwidth} #3 \end{column} ~~~
  \begin{column}{#1\textwidth} #4 \end{column} \end{columns}
}
\newcommand{\frsplitdiffthree}[7][]{
  \begin{columns}[#1]%[T] % align columns
  \begin{column}{#2\textwidth} #5 \end{column} ~~~
  \begin{column}{#3\textwidth} #6 \end{column} ~~~
  \begin{column}{#4\textwidth} #7 \end{column} \end{columns}
}
\newcommand{\frsplittthree}[4][.32]{
  \begin{columns}[T] % align columns
  \begin{column}{#1\textwidth} #2 \end{column} ~
  \begin{column}{#1\textwidth} #3 \end{column} ~
  \begin{column}{#1\textwidth} #4 \end{column} \end{columns}
}


\newcommand{\typerule}[4][]{\ensuremath{\begin{array}[#1]{c}\textsf{\scriptsize ({#2})} \\#3 \\\hline\raisebox{-3pt}{\ensuremath{#4}}\end{array}}}
\newcommand{\styperule}[3][]{\ensuremath{\begin{array}[#1]{c} #2 \\[0.5mm]\hline\raisebox{-4pt}{\ensuremath{#3}}\end{array}}}
\newcommand{\shrk}{\vspace{-3mm}}

\def\caixa#1{\medskip
  \begin{center}
  \fbox{\begin{minipage}{0.9\textwidth}\protect{#1}\end{minipage}}
  \end{center}}

\newcommand{\mybox}[2][4mm]{
  % \begin{minipage}{#1\textwidth}\begin{block}{}\centering #2\end{block}\end{minipage}}
  \tikz{\node[fill=barcolor!60,align=center,inner sep=#1]{#2};}}
\newcommand{\mycbox}[2][4mm]{
  \begin{center}\mybox[#1]{#2}\end{center}}
  % {\\[-5mm]\centering\mybox[#1]{#2}\\[-5mm]}}

%%%%% Tikz
% \usetikzlibrary{arrows.meta, calc, fit, tikzmark}
\usetikzlibrary{%
  positioning
 ,patterns
 ,arrows
 ,arrows.meta
 ,automata
 ,calc
 ,shapes
 ,fit
 ,tikzmark
 ,fadings
 ,decorations.pathreplacing
 ,plotmarks
% ,pgfplots.groupplots
 ,decorations.markings
 ,shadows
}
% \tikzset{shorten >=1pt,node distance=2cm,on grid,auto,initial text={},inner sep=2pt}
\tikzstyle{aut}=[shorten >=1pt,node distance=2cm,on grid,auto,initial text={},inner sep=2pt]
\tikzstyle{st}=[circle,draw=black,fill=black!10,inner sep=3pt]
\tikzstyle{sst}=[rectangle,draw=none,fill=none,inner sep=3pt]
\tikzstyle{final}=[accepting]

%%% Uppaal-like diagrams
\newcommand{\uppbox}[3][20mm]{\tikz{
  \node[black!15,fill=black!15,minimum width=#1,align=left](title){\textbf{{\footnotesize #2}}};
  \node[black!15,fill=black!15,left,xshift=4mm]at(title.east){\textbf{{\footnotesize #2}}};
  \node[blue!60!cyan,right] at(title.west){\textbf{{\footnotesize #2}}};
  \node[below,inner sep=2mm,fill=white,xshift=2mm](box)at(title.south){\includegraphics[width=#1]{#3}};
  \node[fit=(title)(box),draw=black,inner sep=0pt]{};
}}

\newcommand{\uppboxv}[3][20mm]{\tikz{
  \node[below,inner sep=2mm,fill=white](box){\includegraphics[height=#1]{#3}};
  \coordinate[yshift=5mm](top)at(box.north);
  \node[fit=(top)(box.north west)(box.north east),inner sep=0pt,fill=black!15](title){};
  \node[blue!60!cyan,right] at(title.west){\textbf{{\footnotesize #2}}};
  \node[fit=(title)(box),draw=black,inner sep=0pt]{};
}}

%%% frame for pictures [graphics-options]{content}
\newcommand{\includegraphicsframed}[2][]{\tikz{
  \node[below,inner sep=1mm,fill=barcolor,%xshift=2mm,
    draw=black,
    drop shadow={
      top color=gray,
      bottom color=white,
      %fill=gray,
      opacity=0.2,
      shadow xshift=4pt,
      shadow yshift=-3pt
    }](box){\includegraphics[#1]{#2}};
}}


%% COnfiguring Listings
\lstset{ % basic style
  language=scala,
  basicstyle=\ttfamily\scriptsize,
  breakatwhitespace=true,
  breaklines=true,
  % mathescape,
  % morecomment=[l]{//},
  % morecomment=[n]{/*}{*/},
  % frame=single,                    % adds a frame around the code
  rulecolor=\color{black!40},         % if not set, the frame-color may be changed on line-breaks within not-black text (e.g. comments (green here))
  xleftmargin=1.5mm,
  xrightmargin=1.5mm,
  backgroundcolor=\color{black!5},
  % line numbers
%  numbers=left,  % where to put the line-numbers; possible values are (none, left, right)
 numbersep=5pt, % how far the line-numbers are from the code
 numberstyle=\tiny\color{gray},   
 stepnumber=1,  % the step between two line-numbers. If it is 1 each line will be numbered      
%  xleftmargin=3mm,
%  xrightmargin=1.5mm,
%%%%%
  captionpos=b, % t or b (top or bottom)
  belowcaptionskip=5mm,
%%%%%
  % alsoletter={-},
  % emphstyle=\ttfamily\color{blue}, %\underbar,
  % emphstyle={[2]\ttfamily\color{green!50!black}},
  emphstyle=\bfseries\itshape\color{blue!80!black},       % moreemph={...} - layer keywords
  emphstyle={[2]\itshape\color{red!70!black}},%\underbar} % moreemph={[2]...} - inner keywords
  %
  keywordstyle=\bf\ttfamily\color{red!50!black},
  commentstyle=\sl\ttfamily\color{gray!70},
  % commentstyle=\color{green!60!black},
  stringstyle=\ttfamily\color{green!60!black},
  morestring=[b]",
  % morecomment=[l]{\#},
  frame=single,
  % numberstyle=\tiny, numbers=left, stepnumber=1, firstnumber=1, numberfirstline=true,
  %emph={act,proc,init,sort,map,var,eqn},
  %emph={[2]block,hide,comm,rename,allow,||,<>,sum,&&,=>,true,false},
  literate={\\§}{{{\mbox{\textdollar}}}}1,
  % literate=*{->}{{{\color{red!70!black}$\to$}}}{1}
  %            {.}{{{\color{red!70!black}.}}}{1}
  %            {+}{{{\color{red!70!black}\hspace*{1pt}+\hspace*{1pt}}}}{1}
  %            {|}{{{\color{red!70!black}|}}}{1}
  %            {||}{{{\color{red!70!black}|\!\!|}}}{1}
  %            {*}{{{\color{red!70!black}*}}}{1}
  %            {\#}{{{\color{red!70!black}\#}}}{1}
  %            {&}{{{\color{red!70!black}\&}}}{1}
  %            {:=}{{{\color{red!70!black}:=}}}{1}
  %            {=>}{{{\color{red!70!black}=>}}}{1}
  %            % {>}{{{\color{green!65!black}\hspace*{1pt}>\hspace*{1.5pt}}}}{2}
  %            % {<}{{{\color{green!65!black}\hspace*{1.5pt}<\hspace*{1pt}}}}{2}
  %            % {]}{{{\color{green!65!black}\hspace*{1pt}]\hspace*{1.5pt}}}}{1}
  %            % {[}{{{\color{green!65!black}\hspace*{1.5pt}[\hspace*{1pt}}}}{1}
  %            {>}{{{\color{green!65!black}>}}}{1}
  %            {<}{{{\color{green!65!black}<}}}{1}
  %            {]}{{{\color{green!65!black}]}}}{1}
  %            {[}{{{\color{green!65!black}[}}}{1}
  %  morekeywords={Merger1,Fifo2,Lossy3,Init1,Init2}
  % ,emph={act,proc,init,sort,map,var,eqn}
  % ,emph={[2]block,hide,comm,rename,allow,||,<>,sum,&&,=>}
}

\newcommand{\code}[1]{{\relsize{-1}\ttfamily #1}}
\newcommand{\cod}[1]{{\relsize{-0.5}\color{blue!75!black}\ttfamily #1}}
% \definecolor{mydarkgreen}{rgb}{0,0.6,0}
% \newcommand{\bash}[1]{\lstinline[basicstyle=\ttfamily\relsize{-0.5}\color{mydarkgreen},keywordstyle=\bf\sffamily\color{purple},columns=fullflexible,keepspaces,literate=*]|#1|}

\lstdefinestyle{tiny}{basicstyle=\ttfamily\relsize{-7}}

% Include slides from others
\newcommand{\byothers}[3]{{
\begin{frame}{}~\mycbox{\Large slides by #1\\pages #2}\end{frame}{}
\setbeamercolor{background canvas}{bg=}
\includepdf[pages=#2]{../../others/#3}
}}

\newcounter{mypage}
\newcommand{\fromBook}[6][scale=0.5]{
\setcounter{mypage}{20}
\addtocounter{mypage}{#2}
\begin{tabular}{@{}r@{}}
\includegraphics[page=\themypage, trim = #3 #4 #5 #6, clip, #1]%
  {../../learning-concurrent-scala/learning-concurrent-programming-in-scala.pdf}
  \\[0mm]
  \textcolor{myGray}{$\left[\begin{array}{@{}r@{}}
  ~\\[-7mm]
  \text{{\tiny in \emph{``Learning Concurrent}}}\\[-3mm]
  \text{{\tiny \emph{Programming in Scala''}, pg.\,#2}}
  \\[0mm]
  \end{array}\right]$}
\end{tabular}
}

\newcommand{\fromBookW}[4][width=\textwidth]{
\setcounter{mypage}{20}
\addtocounter{mypage}{#2}
\begin{tabular}{@{}r@{}}
\includegraphics[page=\themypage, trim = 20mm #3 20mm #4, clip, #1]%
  {../../learning-concurrent-scala/learning-concurrent-programming-in-scala.pdf}
  \\[-2mm]
  \textcolor{myGray}{$\left[\begin{array}{@{}r@{}}
  ~\\[-7mm]
  \text{{\tiny in \emph{``Learning Concurrent}
                      \emph{Programming in Scala''}, pg.\,#2}}
  \\[0mm]
  \end{array}\right]$}
\end{tabular}
}

\newcommand{\fromBookWframed}[4][width=\textwidth]{
\setcounter{mypage}{20}
\addtocounter{mypage}{#2}
\begin{tabular}{@{}r@{}}
\includegraphicsframed[page=\themypage, trim = 20mm #3 20mm #4, clip, #1]%
  {../../learning-concurrent-scala/learning-concurrent-programming-in-scala.pdf}
  \\[-2mm]
  \textcolor{myGray}{$\left[\begin{array}{@{}r@{}}
  ~\\[-7mm]
  \text{{\tiny in \emph{``Learning Concurrent}
                      \emph{Programming in Scala''}, pg.\,#2}}
  \\[0mm]
  \end{array}\right]$}
\end{tabular}
}

\newcommand{\fromAuthor}[3][width=\textwidth]{
\begin{tabular}{@{}r@{}}
\includegraphics[#1]%
  {src/img/#3}
  \\[-2mm]
  \textcolor{myGray}{$\left[\begin{array}{@{}r@{}}
  ~\\[-7mm]
  \text{{\tiny in #2}}
  \\[0mm]
  \end{array}\right]$}
\end{tabular}
}


%-------------- template --------------------------------------------------
 %!TEX root=../../1-intro.tex
 \usetheme{metropolis}
\usepackage{appendixnumberbeamer}

%------ Setting lecture info ----------------------------------------------
\newcounter{lectureID}
\stepcounter{lectureID}
\newcommand{\getLecture}{\arabic{lectureID}\xspace}
\newcommand{\setLectureBasic}[1]{
  \title{
    #1
    }
  \author{Nelma Moreira ~~\&~~ \textbf{Jos\'{e} Proen\c{c}a}}
  \institute{CISTER -- U.Porto, Porto, Portugal
            \hfill 
            \begin{tabular}{r@{}}
            \url{https://fm-dcc.github.io/cp2425}
            \end{tabular}
            }
  \date{Concurrent programming (CC3040) 2024/2025}
  % logos of institutions
  \titlegraphic{
    \begin{textblock*}{5cm}(2.0cm,7.00cm)
       \includegraphics[scale=0.14]{src/img/logos/fcup}\hspace*{.85cm}~%
    \end{textblock*}
    \begin{textblock*}{5cm}(6.0cm,7.25cm)
      \includegraphics[scale=0.43]{src/img/logos/cmup}
    \end{textblock*}
    \begin{textblock*}{5cm}(10.4cm,7.45cm)
      \includegraphics[scale=0.16]{src/img/logos/cister}
    \end{textblock*}
  }  
  \frame[plain]{\titlepage}
}
\newcommand{\setLecture}[2]{\setcounter{lectureID}{#1}\setLectureBasic{#1. #2}}

%------ Counters for exercises ----------------------------------------------
\newcounter{cExercise}
\newcommand{\exercise}{\stepcounter{cExercise}Ex.\,\arabic{lectureID}.\arabic{cExercise}:\xspace}
\newcommand{\exerciseBack}{\addtocounter{cExercise}{-1}}
\newcommand{\exerciseAdd}{\stepcounter{cExercise}}
\newcommand{\doExercise}[3][0mm]{\begin{exampleblock}{\exercise #2}\wrap{\rule{0pt}{#1}}#3\end{exampleblock}}
\newcommand{\doSimpleExercise}[2][0mm]{\begin{exampleblock}{}\wrap{\rule{0pt}{#1}}\structure{\textbf{\exercise} #2}\end{exampleblock}}

% Slide
\newenvironment{slide}[1]{\begin{frame}\frametitle{#1}}{\end{frame}}



% Base colors (from metropolis theme)
\definecolor{metDarkBrown}{HTML}{604c38}
\definecolor{metDarkTeal}{HTML}{23373b}
\definecolor{metLightBrown}{HTML}{EB811B}
\definecolor{metLightGreen}{HTML}{14B03D}
\definecolor{metDarkGreen}{HTML}{4f9277}

 

\metroset{numbering=fraction,progressbar=frametitle}

% \setbeamercolor*{structure}{fg=blue!80!black}
\setbeamercolor*{structure}{fg=metLightGreen}

% % \definecolor{MainColour}{rgb}{0., 0.25, 0.8}
% \colorlet{MainColour}{blue!50!black}
% \colorlet{BgColour}{blue!10}
% \colorlet{BarColour}{blue!50!black}

% %\usetheme{CambridgeUS}%{Copenhagen}%{Frankfurt}%{Singapore}%{CambridgeUS}
% \usecolortheme[named=MainColour]{structure} 
% \useoutertheme[subsection=false]{miniframes}
% \useinnertheme{circles}
% %\useinnertheme[shadow=false]{rounded}
% \setbeamertemplate{blocks}[rounded][shadow=false]

% \setbeamercovered{transparent} 
% \setbeamertemplate{navigation symbols}{} %Remove navigation bar
% \setbeamertemplate{footline}[frame number] % add page number
% \setbeamercolor{postit}{fg=MainColour,bg=BgColour}
% \setbeamercolor{structure}{bg=black!10}
% %\setbeamercolor{palette primary}{use=structure,fg=red,bg=green}
% %\setbeamercolor{palette secondary}{use=structure,fg=red!75!black,bg=green}
% \setbeamercolor{palette tertiary}{use=structure,bg=BarColour,fg=white}
% %\setbeamercolor{palette quaternary}{fg=black,bg=green}
% %\setbeamercolor{normal text}{fg=black,bg=white}
% %\setbeamercolor{block title alerted}{fg=red,bg=green}
% %\setbeamercolor{block title example}{bg=black!10,fg=green}
\setbeamercolor{block body}{bg=black!5}

% \setbeamercolor{block title alerted}{bg=red!25}
% \setbeamercolor{block body alerted}{bg=red!10}

% \setbeamercolor{block title example}{bg={rgb:green,2;black,1;white,5}}
\setbeamercolor{block body example}{bg={rgb:green,2;black,1;white,20}}
\setbeamercolor{block body alerted}{bg={metLightBrown!25}}

% \setbeamertemplate{itemize item}{\color{black!10}$\blacksquare$}
\setbeamercolor{itemize item}{fg=metDarkTeal}
\setbeamercolor{itemize subitem}{fg=metDarkTeal}

\setbeamercolor{graybc}{fg=black,bg=black!10}
\newcommand{\myblock}[1]{\begin{beamercolorbox}[dp=1ex,center,rounded=true]%
  {graybc} {\large \textbf{#1}} \end{beamercolorbox}}%

%%%%%%%%%

% \definecolor{barcolor}{rgb}{.65,.79,.92} % FCUP color
\colorlet{barcolor}{metDarkGreen!60}


% Configuring the foot line
\setbeamercolor{author in head/foot}{fg=metDarkTeal, bg=barcolor}%
\setbeamercolor{date in head/foot}{fg=barcolor!75!black}%
\setbeamertemplate{footline}
{
  \leavevmode%
  \hbox{%
  \begin{beamercolorbox}[wd=.4\paperwidth,ht=2.25ex,dp=1ex,center]{author in head/foot}%
    \usebeamerfont{author in head/foot}\insertshortauthor
  \end{beamercolorbox}%
  \begin{beamercolorbox}[wd=.5\paperwidth,ht=2.25ex,dp=1ex,center]{title in head/foot}%
    \usebeamerfont{title in head/foot}\insertsection
  \end{beamercolorbox}%
  \begin{beamercolorbox}[wd=.1\paperwidth,ht=2.25ex,dp=1ex,right]{date in head/foot}%
    \insertframenumber{} / \inserttotalframenumber\hspace*{2ex} 
  \end{beamercolorbox}}%
  \vskip0pt%
}
% No configuration symbols
\setbeamertemplate{navigation symbols}{}


%%%%%%% Custom bar above

\setbeamercolor{frametitle}{bg=barcolor,fg=metDarkTeal}
\setbeamercolor{progress bar}{fg=barcolor!75!black}

%%% Custom bar with a LOGO

\makeatletter
\setbeamertemplate{frametitle}{%
  \nointerlineskip%
  \begin{beamercolorbox}[%
      wd=\paperwidth,%
      sep=0pt,%
      leftskip=\metropolis@frametitle@padding,%
      rightskip=\metropolis@frametitle@padding,%
    ]{frametitle}%
  \metropolis@frametitlestrut@start%
  \insertframetitle%
  \nolinebreak%
  \metropolis@frametitlestrut@end%
  \hfill
  \raisebox{-1.3ex}{\includegraphics[height=4ex,keepaspectratio]{src/img/logos/fcup-bar}}
  \end{beamercolorbox}%
}
\makeatother





%----------------------------------------------------------------------------

\begin{document}

% \setLectureBasic{Cyber-Physical Computation}
\setLecture{2}{Short introduction to Scala}
% \frame[plain]{\titlepage}


\section{Motivation}

\begin{frame}\frametitle{Why Scala?}
    
  \splitthree{0.31}{0.31}{0.35}{
    \begin{block}{Used by many modern concurrency frameworks}
      - Syntactic flexibility
      \\- Programming models as Embedded Domain Specific Languages
      \\- Many useful features
    \end{block}
  }{
    \begin{exampleblock}{Safe language}
      - Automatic garbage collection
      \\- Automatic bound checks
      \\- No pointer arithmetic
      \\- Static type safety
    \end{exampleblock}
  }{
    \begin{alertblock}{Java interoperability}
      - Compiled to Java bytecode
      \\- Can use existing Java libraries
      \\- Good interaction with Java's rich ecosystem
      \\- Chosen by some Java-compatible frameworks
    \end{alertblock}
  }

\end{frame}

\section{Executing Scala}

\begin{frame}[fragile]{Running with SBT}
% \centering
~\\[-8mm]
% \doSimpleExercise{What do these functions do?}
\begin{columns}
\begin{column}{0.48\textwidth}
\begin{lstlisting}
object SquareOf5 extends App {
  def square(x: Int): Int = x * x
  val s = square(5)
  println(s"Result: \§s")
}
\end{lstlisting}
\end{column}
\begin{column}{0.48\textwidth}
\structure{Call stack} vs. \alert{object heap}
\\[2mm]
where are values stored?
\\[5mm]
\begin{block}{Concurrent threads:}
do not share the \structure{call stack},
\\share the \alert{object heap},  
\end{block}
% \begin{lstlisting}[language=scala]
% ...
% \end{lstlisting}
\end{column}
\end{columns}

\pause
\bigskip
~~~~~~\alert{Local:} SBT (practical lessons)\\
~~~~~~\alert{Online:} \url{https://scastie.scala-lang.org/cIf3BCTQRBybqMcQpYArGA}
\end{frame}


\begin{frame}\frametitle{Stack and Heap}
    \centering
    
    % \fromBook{19}{45mm}{98mm}{43mm}{35mm} % left,bottom,right,up
    \fromBookW{19}{103mm}{63mm}

\end{frame}

\section{Scala in a nutshell}

\begin{frame}[fragile]\frametitle{Creating a Printer class}
~\\[-8mm]
% \doSimpleExercise{What do these functions do?}
\begin{columns}
\begin{column}{0.52\textwidth}
\begin{lstlisting}
class Printer(val greeting: String) {
  def printMessage(): Unit = println(greeting + "!")
  def printNumber(x: Int): Unit = {
    println("Number: " + x)
  }
}
\end{lstlisting}
\end{column}
\begin{column}{0.44\textwidth}
\only<1>{One \texttt{greeting} \alert{parameter}
\\Two \alert{methods}}%
\pause
% \only<2>{
\structure{\textbf{Using the class}}
\begin{lstlisting}
val printy = new Printer("Hi")
          // instantiate
printy.printMessage() // ?
printy.printNumber(5) // ?
\end{lstlisting}
% }
\end{column}
\end{columns}
\end{frame}


\begin{frame}[fragile]\frametitle{Singleton objects}
~\\[-8mm]
\begin{columns}
\begin{column}{0.44\textwidth}
\begin{lstlisting}
object Test {
  val Pi = 3.14
}
\end{lstlisting}
\end{column}
\begin{column}{0.52\textwidth}
\structure{\textbf{Using the object}}
\begin{lstlisting}
val x = Test.Pi * 5 * 5
    // no need to instantiate
\end{lstlisting}% }
\end{column}
\end{columns}
\end{frame}


\begin{frame}[fragile]\frametitle{Traits (similar to Java Interfaces)}    
~\\[-8mm]
\begin{columns}
\begin{column}{0.58\textwidth}
\begin{lstlisting}
trait Logging {
  def log(s: String): Unit // just declared
  def warn(s: String) = log("WARN: " + s)
  def error(s: String) = log("ERROR: " + s)
}
class PrintLogging extends Logging {
  def log(s: String) = println(s)
}
\end{lstlisting}
\end{column}
\begin{column}{0.40\textwidth}
\structure{\textbf{Using traits}}
\begin{lstlisting}
val x = new PrintLogging
val y = new Logging {
  def log(s:String): Unit =
    println(s)
}
\end{lstlisting}% }
\end{column}
\end{columns}
\end{frame}


\begin{frame}[fragile]\frametitle{Type parameters}    
~\\[-8mm]
\begin{columns}
\begin{column}{0.58\textwidth}
\begin{lstlisting}
class Pair[A,B](val fst: A, val snd: B)
\end{lstlisting}
\end{column}
\begin{column}{0.40\textwidth}
\structure{\textbf{Using Pair}}
\begin{lstlisting}
val x: Pair[Int,String] = new Pair(4,"a")
val y = new Pair(2,5) // infer type
\end{lstlisting}% }
\end{column}
\end{columns}
\end{frame}


\begin{frame}[fragile]\frametitle{Lambdas (anonymous functions)}    
~\\[-8mm]
\begin{columns}
\begin{column}{0.58\textwidth}
\begin{lstlisting}
val twice_a: Int=>Int = (x:Int) => x*2
val twice_b = (x:Int) => x*2
val twice_c: Int=>Int = x => x*2
val twice_d: Int=>Int = _ * 2
\end{lstlisting}
\end{column}
\begin{column}{0.40\textwidth}
\structure{\textbf{Using lambdas}}
\begin{lstlisting}
val x = twice_a(4)
\end{lstlisting}% }
\end{column}
\end{columns}
\end{frame}


\begin{frame}[fragile]\frametitle{Byname parameters (lazy)}
~\\[-8mm]
\begin{columns}
\begin{column}{0.49\textwidth}
\begin{lstlisting}
def runTwice(body: =>Unit) = {
  body
  body
}
\end{lstlisting}
\end{column}
\begin{column}{0.49\textwidth}
\structure{\textbf{Using Byname}}
\begin{lstlisting}
runTwice { // prints "Hello" twice
  println("Hello")
}
\end{lstlisting}% }
\end{column}
\end{columns}
\end{frame}


\begin{frame}[fragile]\frametitle{``\texttt{for}'' expressions and comprehension}
~\\[-8mm]
\begin{columns}
\begin{column}{0.50\textwidth}
\begin{lstlisting}[emph={foreach,map,flatMap}]
for (i <- 0 until 10) println(i)
// equivalent to
0.until(10).foreach(i => println(i))

val negatives_a =
  for (i <- 0 until 10) yield -i
val negatives_b =
  (0 until 10).map(i => -1 * i)
\end{lstlisting}
\end{column}
\begin{column}{0.48\textwidth}
\pause
% \structure{``\texttt{for-yield}'' expressions}
\begin{lstlisting}[emph={foreach,map,flatMap}]
val pairs_a =
  for (x <- 0 until 4;
       y <- 0 until 4) yield (x, y)
val pairs_b =
  (0 until 4).flatMap(x =>
    (0 until 4).map(y =>
      (x, y)))
\end{lstlisting}
\end{column}
\end{columns}
\end{frame}


\begin{frame}[fragile]\frametitle{Scala collections and string interpolation}
~\\[-8mm]
\begin{columns}
\begin{column}{0.46\textwidth}
Common collections:
\\~~\texttt{Seq[T]}, \texttt{List[T]}, \texttt{Set[T]}, \texttt{Map[K,V]}  
\begin{lstlisting}
val msgs_a: Seq[String] = Seq("Hello", "world!")
val msgs_b: List[String]= "Hello"::"World"::Nil
val msgs_c: Set[String]= Set("Hello", "world!")
val msgs_d: Map[String,Int]= Map("Hello"->5, "world!"->6)
\end{lstlisting}
\end{column}
\begin{column}{0.52\textwidth}
String interpolation:
\begin{lstlisting}
val number = 7
val msg =
  s"After \§number comes \§{number+1}!"
\end{lstlisting}
\end{column}
\end{columns}
\end{frame}


\begin{frame}[fragile]\frametitle{Pattern matching}
~\\[-8mm]
\begin{columns}
\begin{column}{0.56\textwidth}
\begin{lstlisting}
val successors =
  Map(1 -> 2, 2 -> 3, 3 -> 4)
successors.get(5) match {
  case Some(n) =>
    println(s"Successor is: \§n")
  case None    =>
    println("Could not find successor.")
}
\end{lstlisting}
\end{column}
\begin{column}{0.42\textwidth}
\begin{lstlisting}
trait IntOrError
case class MyInt(i:Int)      extends IntOrError
case class MyError(e:String) extends IntOrError

//...
def show(ie: IntOrError)
  ie match {
    case MyInt(i)   => println(s"Number: \§i")
    case MyError(e) => println(s"Error: \§e")
  }
}
\end{lstlisting}
\end{column}
\end{columns}
\end{frame}


\begin{frame}[fragile]\frametitle{Try-catch}
~\\[-8mm]
\begin{columns}
\begin{column}{0.56\textwidth}
\begin{lstlisting}
trait IntOrError
case class MyInt(i:Int)      extends IntOrError
case class MyError(e:String) extends IntOrError

//...
def getInt(ie: IntOrError)
  ie match {
    case MyInt(i)   => i
    //case MyError(e) => ???
    // (match error if an error is found)
  }
}
\end{lstlisting}
\end{column}
\begin{column}{0.42\textwidth}
\begin{lstlisting}
def showInt(ie:IntOrError) = 
  try {
    println(s"Int value: \§{getInt(i)}")
  } catch { 
    case _: Throwable => println("Got some error - probably not an int.")
  } finally {
    // always executes
    // (for side-effects, not to return results)
    println("Done showing")
  }
\end{lstlisting}
\end{column}
\end{columns}
\end{frame}

\begin{frame}[fragile]\frametitle{Operator overloading}
~\\[-8mm]
\begin{columns}
\begin{column}{0.56\textwidth}
\begin{lstlisting}
class Position(val x: Int, val y: Int) {
  def +(that: Position) =
    new Position(x + that.x, y + that.y)
  def *(n: Int) =
    new Position(x * n, y * n)
}
\end{lstlisting}
\end{column}
\begin{column}{0.42\textwidth}
\structure{Using operators}
\begin{lstlisting}
val p1 = new Position(3,4)
val p2 = p1 + p1 * 2 //?
\end{lstlisting}
\end{column}
\end{columns}
\end{frame}


\begin{frame}[fragile]\frametitle{Package objects}
~\\[-8mm]
\begin{columns}
\begin{column}{0.62\textwidth}
\alert{\textbf{File \texttt{src/main/scala/cp/lablessons/package.scala}}}
\begin{lstlisting}[emph={log}]
package cp

package object lablessons {
  def log(msg: String): Unit =
    println(
      s"\§{Thread.currentThread.getName}: \§msg"
    )
}
\end{lstlisting}
\end{column}
\begin{column}{0.40\textwidth}
The \texttt{log} function is used throughout these lessons
\\[4mm]
Requires starting with\\~~\texttt{package cp.lablessons}
\end{column}
\end{columns}
\end{frame}


\begin{frame}\frametitle{Wrap up}
~\\[-8mm]
\begin{columns}
\begin{column}{0.48\textwidth}
\begin{itemize}
  \item Stack and Heap
  \item Singleton objects
  \item Traits (similar to Java Interfaces)
  \item Type parameters
  \item Lambdas (anonymous functions)
  \item Byname parameters (lazy)
\end{itemize}
\end{column}
\begin{column}{0.54\textwidth}
\begin{itemize}
  \item ``\texttt{for}'' expressions
  \item ``\texttt{for}'' comprehension
  \item Scala collections and string interpolation
  \item Pattern matching
  \item Try-catch
  \item Operator overloading
  \item Package objects
\end{itemize}
\end{column}
\end{columns}

    


\end{frame}

\end{document}
